\section{Purpose of the Supporting Information}

This supplement records implementation details that should not carry the main article's narrative burden. It makes the model schedule, scenario catalog, validation artifacts, and source-panel gaps easy to inspect.

\section{Model Schedule and State}

Each simulation tick selects one policy contest and resolves it through the active institutional arena. The state update order is fixed:

\begin{enumerate}
  \item Select a contest from the scenario contest mix using the seeded random source.
  \item Replenish lobby budgets through client funding when expected private gain, exposure, secrecy needs, or reform threat justify additional money.
  \item Evaluate substitution pressure for each lobby organization and choose a channel strategy.
  \item Allocate finite spending across direct access, agenda access, information distortion, public campaigns, litigation threats, campaign finance, dark money, revolving-door routes, intermediaries, and defensive reform spending.
  \item Apply channel effects to public support, public benefit, private gain, dark-money influence, campaign finance pressure, litigation threat, revolving-door influence, comment-record distortion, procurement pressure, and hidden influence.
  \item Generate influence-network paths from funders and lobby organizations through channels, intermediaries, venues, and targets.
  \item Triage rulemaking comments into unique information, organized technical content, duplicate templates, and synthetic or low-authenticity comments.
  \item Resolve the contest in its arena under the reform regime, public backlash, hidden substitution pressure, and stochastic noise.
  \item Apply detection and sanctions.
  \item Update lobby strategy memory, issue-specific funding multipliers, regulator queue pressure, watchdog focus, adaptation speed, and reform-decay pressure.
\end{enumerate}

The full ODD-style description is maintained in \texttt{docs/odd-model.md} and is included with the submission bundle as supporting information.

\section{Scenario Families}

The scenario catalog is organized around mechanism families rather than institutional venues alone. The families are:

\begin{itemize}
  \item baseline and disclosed lobbying;
  \item low-salience technical rulemaking;
  \item campaign-finance, dark-money, and outside-spending pressure;
  \item revolving-door and advisory substitution;
  \item think-tank, association, and sponsored-expert intermediaries;
  \item procurement and vendor venue-shifting;
  \item hard lobbying budgets and shadow-lobbying stress;
  \item public financing, democracy vouchers, public-interest representation funds, and public advocate offices;
  \item real-time disclosure, machine-readable meeting logs, beneficial-owner disclosure, comment-authenticity systems, cooling-off enforcement, randomized audits, sanctions, procurement firewalls, and venue-shifting detection;
  \item reform bundles, ablations, interactions, sensitivity sweeps, and portfolio screens.
\end{itemize}

The generated scenario catalog in \texttt{docs/scenario-catalog.md} remains the authoritative scenario reference.

\section{Primary Evaluation Rule}

The primary reform criterion is total influence distortion. A reform is treated as a distortion failure when observed capture falls and total modeled distortion rises. It is treated as a substitution warning when observed capture falls but one or more hidden diagnostics rises:

\begin{align*}
\Delta observedCapture &< 0,\\
\Delta hiddenInfluence &> 0 \quad \text{or}\\
\Delta hiddenCapture &> 0 \quad \text{or}\\
\Delta totalDistortion &> 0 \quad \text{or}\\
\Delta substitutionRisk &> 0.
\end{align*}

This diagnostic is implemented in \texttt{scripts/validate-reports.py}. The generated outputs are \texttt{reports/substitution-audit.csv}, \texttt{reports/substitution-audit.md}, plus the substitution-warning table and map in the main manuscript.

\section{Generated Report Artifacts}

The reproducible report pipeline writes:

\begin{itemize}
  \item campaign, mechanism-comparison, sensitivity, ablation, interaction, and portfolio CSV and Markdown reports;
  \item source moments from normalized public-data panels;
  \item source-panel inventory diagnostics;
  \item validation summary and classified calibration queue;
  \item substitution audit with distortion-failure, warning, and channel-shift tradeoff labels;
  \item generated paper tables, figures, local PDF, Wiley PDF, supplement PDF, and Wiley submission archive.
\end{itemize}

The canonical local check is:

\begin{verbatim}
make paper-artifacts-check
\end{verbatim}

\section{Validation Layers}

Validation proceeds in four layers. Parser and fixture tests enforce source schemas before normalized source rows enter calibration. The source-moments target extracts direct and proxy distributional moments for lobbying, campaign finance, rulemaking, procurement, intermediaries, public financing, and revolving-door panels. The report-validation target compares model outputs against broad benchmark ranges and the parameter map, while the substitution audit distinguishes clean improvements, distortion failures, and substitution warnings. Channel movement without higher hidden influence or distortion is treated as a tradeoff rather than a failure.

The portfolio target ranks candidate reform bundles by total distortion, hidden capture, substitution risk, administrative burden, network opacity, legitimate-advocacy chill, speech-restriction risk, cross-venue detection, and participation protection. The calibration-queue target classifies misses into model tuning, metric splits, direct source moments, scenario coverage, scale alignment, and benchmark review. Future snapshots can add health, transportation, broader SAM/FPDS procurement panels, DIME, LobbyView, IRS 527 and Form 990 rows, FACA and witness records, Seattle voucher rows, direct electioneering or dark-money source exports, and richer revolving-door source exports once the environmental slice is stable.

\section{Source-Panel Coverage}

The current 2024 EPA/ENV snapshot is a reproducibility scaffold, not a representative empirical panel. It is strongest for LDA, OpenFEC Schedule E, Regulations.gov, Federal Register, and USAspending schema checks. Public financing and direct dark money are separate bridge panels. This prevents a public-financing fixture or future dark-money export from being mistaken for ordinary FEC receipts, and it keeps direct dark money separate from super PAC outside spending. The source bridge remains incomplete for direct dark-money source rows, representative public-financing program data, broad SAM/FPDS procurement panels, representative nonprofit/intermediary routing, and post-employment movement beyond LDA-derived covered-position records.

The source-panel inventory in \texttt{reports/source-panel-inventory.md} is included in the submission bundle. It should be read before using any model output as more than a mechanism diagnostic.

\section{Decision Rules and Parameters}

The implementation exposes mechanism-facing parameters rather than estimated coefficients. The most important rules are:

\begin{itemize}
  \item Client funding increases with issue preference, exposure, private gain, calibration issue scale, secrecy incentives, and reform-threat pressure. Public-match and voucher sources lower large-donor dependence; dark-money sources increase opacity and legal-risk exposure.
  \item Donor concentration in reports now blends the simulated contribution ledger with source-derived top-donor concentration, Herfindahl concentration, and large-donor shares from the campaign-finance bridge. This avoids treating a report-level donor-Gini proxy as the only donor-network anchor.
  \item Procurement bridge risk now separates initial awards from post-award modifications and uses recipient concentration, single-bid or limited-competition indicators, modification share, price-only awards, protest flags, exclusion flags, and firewall coverage.
  \item Substitution activates when visible-channel pressure, anonymity value, messenger credibility, intermediary fit, venue overlap, and evasion freedom exceed legal cost, disclosure cost, cross-venue detection, and setup time. The threshold is deliberately low enough for stress scenarios to reveal channel migration.
  \item A distortion failure requires more than channel movement. The audit flags a row as a failure when observed or visible capture falls while total distortion rises, and as a warning when hidden influence, hidden capture, or substitution risk rises.
\end{itemize}

The source-facing parameters and benchmark ranges are maintained in \path{data/calibration/parameter-map.csv}. Report-level bounded diagnostics remain synthetic unless the validation summary classifies the target as a source moment.

\section{Figure and Layout Audit}

Paper figures are generated as SVGs and converted to Wiley-preferred PDF graphics. The script \texttt{scripts/generate-interaction-figures.py} uses simple label-placement logic with visible callout lines so point labels remain distinct and connected to the corresponding point. The script \texttt{scripts/audit-paper-layout.py} checks the generated PDFs for sparse float pages, excessive vertical whitespace, and figure or table pages with too little surrounding readable text. Its report is written to \texttt{reports/paper-layout-audit.md}. A separate visual checklist is written to \texttt{reports/manual-visual-audit.md}; it checks generated SVG label boxes for overlap, bounds, and leader-line coverage where callout labels exist and records table-float layout status.

\section{Remaining Calibration Priorities}

The main paper should not claim calibrated policy estimates until the following source panels and tuning loops improve:

\begin{itemize}
  \item direct dark-money identifiers and intermediary routing panels;
  \item representative public-financing and voucher participation records;
  \item broader SAM/FPDS procurement identifiers, modifications, offer counts, exclusions, protests, and firewalls;
  \item richer revolving-door and advisory-access records;
  \item comment-authenticity, duplicate-compression, detection, sanction, and regulator-backlog calibration;
  \item scenario-specific benchmarks for high-end outside spending, hidden substitution, and cooling-off evasion.
\end{itemize}

The calibration queue in \texttt{reports/calibration-queue.md} records the classified work list.

\section{How to Read the Bundle}

The supporting files are organized to make the boundary between model assumptions and evidence visible. The source moments and source-panel inventory should be read first when assessing empirical scope. They state which panels are direct, proxy, thin, fixture-supported, or missing. The mechanism-comparison report should be read next because it shows the consequence of enabling substitution relative to simpler model modes. The campaign, sensitivity, ablation, interaction, and portfolio reports are then best interpreted as stress tests under those model assumptions.

For replication, the Makefile targets regenerate the committed reports, tables, figures, local PDF, Wiley PDF, supplement, layout audit, visual checklist, and submission ZIP from the same checked-in sources. Credentials and ignored raw API payloads are not needed for this default rebuild because the submission path uses normalized committed snapshots. Live-source refreshes are intentionally separate so exploratory API output cannot silently alter the review bundle.
